%%%%%%%%%%%%%%%%%%%%%%%%%%%%%%%%%%%%%%%%%%%%%%%%%%%%%%%%%%%%%%%%%%
%% This is the `x_poluni.tex' included LaTeX file.              %%
%% Toto je inkludovaný LaTeXový soubor `x_poluni.tex'.          %%
%% ------------------------------------------------------------ %%
%% Verze / Version 2022-04-30                                   %%
%% Developed & maintained by / Vyvíjí & udr¾uje                 %%
%%                           <stanislav.hledik@physics.slu.cz>  %%
%%%%%%%%%%%%%%%%%%%%%%%%%%%%%%%%%%%%%%%%%%%%%%%%%%%%%%%%%%%%%%%%%%

%###############################################################################
\chapter{Configurations of uniform density}\label{Chap:Uniform}

It is interesting that there is a special class of general relativistic
polytropes of the index $n = 0$ where the structure equations can be
integrated in terms of elementary functions. We shall discuss these polytropes
in detail because they can give an intuitive insight into the role of the
cosmological constant and can serve as a test bed for the general case of
polytropes with $n>0$.

%%%%%%%%%%%%%%%%%%%%%%%%%%%%%%%%%%%%%%%%%%%%%%%%%%%%%%%%%%%%%%%%%%%%%%%%%%%%%%%%
\section{Structure of the polytropes}\label{Sec:UniStr}

The $n=0$ polytropes correspond to the special class of the internal \SaodS{}
spacetimes~\cite{Stu:2000:ACTPS2:} where the distribution of density $\rho$ is
uniform although the pressure grows monotonically from its zero value on the
surface of the configuration to a maximum value at its center. Recall that in
the configurations with $\rho = \mathrm {const}$, it is not necessary to use
the unrealistic notion of an incompressible fluid~-- one can consider fluids
with pressure growing as radius decreases, being ``hand
tailored''~\cite{Mis-Tho-Whe:1973:Gra:}.  Assuming $n=0$, equations of
structure can be integrated yielding
\begin{equation}
    v(\xi) = \frac{1}{3}\,\xi^{3}\,,                       \label{Eq:grp87}
\end{equation}
and
\begin{equation}
  \left[3-2\sigma(1+\lambda)\xi^{2}\right]\,\oder{\theta}{\xi}
  + (1-2\lambda+3\sigma\theta)(1+\sigma\theta) = 0         \label{Eq:grp88}
\end{equation}
and can be integrated directly after separation of variables. Using the
boundary condition $\theta(0)=1$, we obtain
\begin{equation}
  \sigma\theta = \frac%
    {(1-2\lambda+3\sigma)%
      \left[1-\frac{2}{3}\,\sigma(1+\lambda)\xi^{2}\right]^{1/2}
      -(1-2\lambda)(1+\sigma)}%
    {3(1+\sigma)-(1-2\lambda+3\sigma)
      \left[1-\frac{2}{3}\,
      \sigma(1+\lambda)\xi^{2}\right]^{1/2}}\,.            \label{Eq:grp89}
\end{equation}
This solution determines the dependence of pressure on the radial coordinate,
since for $n=0$ there is $\theta = p(r)/\pcent$. The dependence is
given in units of the energy density since $\sigma =
\pcent/\rhocent$. From the condition $\theta(\xi_{1})= 0$, we find the
radius of the configuration to be determined by
\begin{equation}
  \xi_{1}^{2} = \frac{6[1+2\sigma-\lambda(2+\sigma)]}%
                 {(1-2\lambda+3\sigma)^{2}}\,.             \label{Eq:grp90}
\end{equation}
It is convenient to express the pressure profile in terms of $r, R$, instead
of $\xi, \xi_{1}$, and $\sigma$ obtaining thus the form of expression of the
$n=0$ polytrope as was discussed in~\cite{Stu:2000:ACTPS2:}. Introducing a new
parameter $a$, having dimension of length, by the relation
\begin{equation}
  \frac{1}{a^{2}}
    = \frac{8\pi}{3}\,(\rhocent + \rhovac)     \label{Eq:grp91}
\end{equation}
we find the relation of $R$, $\sigma$, and $\lambda$ to be given by
\begin{equation}
  \left(1-\frac{R^{2}}{a^{2}}\right)^{\!\!1/2}
  = \frac{(1-2\lambda)(1+\sigma)}{1-2\lambda+3\sigma}\,.      \label{Eq:grp93}
\end{equation}
Then the central pressure $\pcent$ and the pressure profile $p(r)$ can
be expressed in the known
form~\cite{Stu:2000:ACTPS2:,%
  Boh:2004:GENRG2:11Solutions,%
  Boh-Fod:2008:PHYSR4:PerFluidCC}
\begin{align}
  \pcent &= \rhocent\,\frac%
    {(1-2\lambda)\left[1-\left(1-\ffrac{R^{2}}{a^{2}}\right)^{1/2}\right]}%
    {3\left(1-\ffrac{R^{2}}{a^{2}}\right)^{1/2}
      -(1-2\lambda)}\,,                                    \label{Eq:grp94}\\
  p(r) &= \rhocent\,\frac%
    {(1-2\lambda)
      \left[
        \left(1-\ffrac{r^{2}}{a^{2}}\right)^{1/2}
        -\left(1-\ffrac{R^{2}}{a^{2}}\right)^{1/2}
      \right]}%
    {3\left(1-\ffrac{R^{2}}{a^{2}}\right)^{1/2}
      -(1-2\lambda)
      \left(1-\ffrac{r^{2}}{a^{2}}\right)^{1/2}}\,.        \label{Eq:grp95}
\end{align}
The radial metric coefficient is given by the relation
\begin{equation}
  \eto{2\Psi(r)}
    = \left(1-\frac{r^{2}}{a^{2}}\right)^{-1}\,,           \label{Eq:grp96}
\end{equation}
and the total mass reads 
\begin{equation}
  M = \frac{1}{3}\,\Lngt\sigma\xi_{1}^{3}
    = \frac{4\pi}{3}\,\rhocent r^{3}\,. \label{Eq:grp97}
\end{equation}
The temporal metric coefficient is determined by the relations  
\begin{equation}
  \eto{\Phi(r)} = \frac{[1 - \frac{2}{3}\,\sigma(1+\lambda)\xi^{2}]^{1/2}}%
                  {1 + \sigma\theta}                       \label{Eq:grp99}
\end{equation}
and
\begin{equation}
  \eto{\Phi(r)} = \frac%
    {3\left(1-\ffrac{R^{2}}{a^{2}}\right)^{1/2}
      -(1-2\lambda)\left(1-\ffrac{r^{2}}{a^{2}}\right)^{1/2}}%
    {2(1+\lambda)}\,.                                      \label{Eq:grp98}
\end{equation}

The special case of $\lambda = -1$ has to be treated separately as
$1/a^{2}=0$. Now, Eq.\,(\ref{Eq:grp88}) reduces to
\begin{equation}
  \frac{\dif\theta}{(1+\sigma\theta)^{2}}
    = -\xi\,\dif\xi\,,                                    \label{Eq:grp100}
\end{equation}
which leads, after integration with the boundary condition $\theta(0)=1$, to
the formula
\begin{equation}
  \theta = \frac%
    {1-\frac{1}{2}(1+\sigma)\xi^{2}}%
    {1+\frac{1}{2}(1+\sigma)\xi^{2}}\,.                   \label{Eq:grp101}
\end{equation}
The boundary of the configuration is at
\begin{equation}
  \xi_{1}^{2} = \frac{2}{1+\sigma}\,.                     \label{Eq:grp102}
\end{equation}
The central pressure and the pressure profile can be expressed in terms of the
radial coordinates $r$, $R$ in the form
\begin{align}
  \pcent &=
    \frac{\rhocent}{\ffrac{2R}{3M}-1}\,,                  \label{Eq:grp103}\\
  p(r) &=
     \pcent\,\frac{1-\ffrac{r^{2}}{R^{2}}}%
     {1+\frac{1}{\ffrac{2R}{3M}-1}\frac{r^{2}}{R^{2}}}\,.  \label{Eq:grp104}
\end{align}
The metric coefficients have a special form, too. The radial $g_{rr}$
component corresponds to the flat $t = \mathrm{const}$ sections
\begin{equation}
  \eto{2\Psi(r)} = 1\,,                                   \label{Eq:grp105}
\end{equation}
while the temporal $g_{tt}$ component takes the form~\cite{Stu:2000:ACTPS2:}
\begin{equation}
  \eto{\Phi(r)} = 1
    + \frac{3M}{2R}\left(\frac{r^{2}}{R^{2}}-1\right)\,.  \label{Eq:grp107}
\end{equation}
Properties of the $n=0$ general relativistic polytropes were discussed
in~\cite{Stu:2000:ACTPS2:,%
  Stu-etal:2001:PHYSR4:}. Here we concentrate on the
existence of these polytropes in dependence on the cosmological constant.

%%%%%%%%%%%%%%%%%%%%%%%%%%%%%%%%%%%%%%%%%%%%%%%%%%%%%%%%%%%%%%%%%%%%%%%%%%%%%%%%
\section[Existence conditions of the $n=0$ polytropes]%
{Existence conditions and compactness of \mbox{n\,=\,0}
  polytropes}\label{Sec:UniCond}

The reality conditions on the general solution, given by Eqs~(\ref{Eq:grp95}),
(\ref{Eq:grp96}), and (\ref{Eq:grp98}), must guarantee that the pressure is
positive and nondivergent, and the metric coefficients have to be regular at
$r \leq R$.  Therefore, two conditions have to be satisfied:
\begin{equation}
  1-2\lambda > 0                                          \label{Eq:grp108}
\end{equation}
and
\begin{equation}
  3\left(1-\frac{R^{2}}{a^{2}}\right)^{\!\!1/2}
   - (1-2\lambda) > 0\,.                                  \label{Eq:grp109}
\end{equation}
Considering the case of $\lambda = 1/2$, we find
\begin{equation}
  \theta(\xi) = \frac{(1-\sigma\xi^{2})^{1/2}}%
                {1+\sigma + \sigma(1-\sigma\xi)^{1/2}}\,. \label{Eq:grp110}
\end{equation}
The boundary of such a polytrope configuration is at
\begin{equation}
  \xi_{1}^{2} = \frac{1}{\sigma}\,,                       \label{Eq:grp111}
\end{equation}
and it can be shown that
\begin{equation}
  R^{2} = a^{2}\,.                                        \label{Eq:grp112}
\end{equation}
Therefore, the metric coefficients are singular because there is
\begin{equation}
  \eto{-2\Psi(R)} = \eto{2\Phi(R)} = \eto{2\Phi(r)}
                 = 1-\frac{R^{2}}{a^{2}}\,.               \label{Eq:grp113}
\end{equation}
For polytrope configurations of a given $M$ and $\Lambda$, condition
(\ref{Eq:grp108}) gives an upper limit on admissible values of the external
radius $R$. Configurations with \mbox{$\lambda \sim 1/2$} can be considered as
nearly ``geodetical'' since the pressure gradient almost vanishes on the
surface, which is close to the static radius of the external geometry. (For
$\lambda = 1/2$, the surface of the static configuration has to be located at
the horizon of the external spacetime, with $R=\ffrac{3r_{\mathrm{g}}}{2}$;
however, no static configuration can have its boundary at a \mbox{black-hole}
horizon, and such configurations are forbidden.)

The lower limit on the external radius of the $n=0$ polytropes is determined
by the condition (\ref{Eq:grp109}) that can be transformed into the relations
\begin{equation}
  \frac{R^{2}}{a^{2}} < \frac{4(1+\lambda)(2-\lambda)}{9} \label{Eq:grp114}
\end{equation}
and
\begin{equation}
  R > 2M\frac{9}{4(2-\lambda)}\,.                         \label{Eq:grp115}
\end{equation}
For $\lambda = 0$, we obtain the well known limit $R>
(\ffrac{9G}{4c^{2}})M$. The restrictions on physically realistic $n=0$
polytropes can also be transformed into a form containing dimensionless
quantities $x \equiv R/M$, $y \equiv \Lambda M^{2}/3$ (see
Refs~\cite{Stu:2000:ACTPS2:,%
  Stu-etal:2001:PHYSR4:}).

Compactness of the polytropic spheres of the uniform density is given by the
relation
\begin{equation}
  \mathcal{C}(\sigma, \lambda)= \frac{2 \sigma [1- 2 \lambda + \sigma
    (2- \lambda)]}{(1- 2 \lambda + 3 \sigma)^{2}}         \label{Eq:grp-c1}
\end{equation}
that is reduced for $\lambda = 0$ to the formula
\begin{equation}
  \mathcal{C}(\sigma)
    = \frac{2 \sigma (1+ 2 \sigma)}{(1+ 3 \sigma)^{2}}\,. \label{Eq:grp-c2}
\end{equation}
The extremely compact configurations with $\mathcal{C} > 1/3$ can exist, if
the central parameter satisfies the relation
\begin{equation}
  \sigma^{2} \geq \sigma^{2}_{\mathrm{ext}} \equiv
    \frac{2}{3}\lambda\,.                                 \label{Eq:grp-c3}
\end{equation}
As an extremely compact spherical configurations are denoted those having
their surface located under the photon circular orbit of the external
spacetime~\cite{Stu:2008:0803.2530:SpSyStCoUniDens}. It can be
shown~\cite{Stu-etal:2001:PHYSR4:} that in extremely compact configurations
(with $R < 3M$), a stable circular null geodesic exists around which null
geodesics captured by the strong gravitational field are concentrated.
Neutrinos, moving along these bound null geodesics, can influence cooling of
extremely compact neutron stars.  The potential well of the captured geodesics
becomes deeper with the repulsive cosmological constant increasing, while it
gets flatter with the attractive cosmological constant
decreasing~\cite{Stu-etal:2012:GENRG2:NeutrinoTrap}.

%%%%%%%%%%%%%%%%%%%%%%%%%%%%%%%%%%%%%%%%%%%%%%%%%%%%%%%%%%%%%%%%%%%%%%%%%%%%%%%%
\section[Gravitational binding of $n=0$ polytropes]%
{Gravitational binding of \mbox{n\,=\,0} polytropes}%
\label{Sec:UniGraB}

Finally, we give the gravitational energy of the $n=0$ polytropes and their
gravitational binding factor $g$. The total energy
\begin{equation}
  E = M = \frac{4\pi}{3}\,\rhocent r^{3}\,.               \label{Eq:grp116}
\end{equation}
The formula for the total proper energy reads
\begin{equation}
  E_{0} = 4\pi\rhocent\int_{0}^R
    \left(1-\frac{r^{2}}{a^{2}}\right)^{-1/2}\,
     r^{2}\,\dif r\,,                                     \label{Eq:grp117}
\end{equation}
and after integration we obtain 
\begin{equation}
  E_{0} = \frac{3}{2}\,M\left(\frac{a}{R}\right)^{3}
    \left[
      \arcsin\left(\frac{R}{a}\right)
      -\frac{R}{a}\left(1-\frac{R^{2}}{a^{2}}\right)^{1/2}
    \right]\,,                                            \label{Eq:grp118}
\end{equation}
where
\begin{equation}
  \frac{R}{a} = \frac{2(1+\lambda)^{1/2}}{1-2\lambda+3\sigma}
    \left\{\sigma[1+2\sigma
    -\lambda(2+\sigma)]\right\}^{1/2}\,.                  \label{Eq:grp119}
\end{equation}
The gravitational potential energy
\begin{equation}
  \mathcal{G} = M
    \left\{
      1 - \frac{3}{2}\left(\frac{a}{R}\right)^{3}
      \left[
        \arcsin\left(\frac{R}{a}\right)
        -\frac{R}{a}\left(1-\frac{R^{2}}{a^{2}}\right)^{1/2}
      \right]
    \right\}                                              \label{Eq:grp120}
\end{equation}
is naturally negative and the gravitational binding factor $g=\mathcal{G}/E$
is given by
\begin{equation}
  g \equiv \frac{\mathcal{G}}{v} = - \frac{3}{2}\left(\frac{a}{R}\right)^{3}
      \left[
        \arcsin\left(\frac{R}{a}\right)
        -\frac{R}{a}\left(1-\frac{R^{2}}{a^{2}}\right)^{1/2}
      \right]\,.                                          \label{Eq:grp121}
\end{equation}
%
\begin{figure}[p]
\begin{center}
\includegraphics[width=\linewidth]{UniSigGBEn}
\par\vspace{0.8\baselineskip}\par
\includegraphics[width=\linewidth]{UniSigGBEnZ}
\end{center}
\caption[Dependence of binding factor on $\sigma$ for uniform density]%
{\label{UniSigGBEn}Dependence of the gravitational binding
  factor~(\ref{Eq:grp121}) for configurations of uniform density on relativity
  parameter $\sigma$ for cosmological parameter
  $\lambda = 0, 10^{-2}, 5\times10^{-2}, 10^{-1}$, respectively. The ``1''
  index emphasizes the energies are related to the whole
  configuration. \emph{Top:}~Wide range of relativity
  parameter. \emph{Bottom:}~Zoom of the shaded region in the top plot.}
\end{figure}

The gravitational potential energy of the $n=0$ polytropes is represented in
Fig.\,\ref{UniSigGBEn}. The role of the cosmological constant is illustrated
by the sequence of lines constructed for appropriately chosen values of the
cosmological parameter~$\lambda$.

%
It expresses the gravitational binding in a ``pure'' form as there in no
internal energy in ``incompressible'' configurations. On the other hand, we
can directly conclude that the relation of the binding energy of gas
$\mathcal{B}$, and the gravitational energy are given by
\begin{equation}
  \frac{\mathcal{B}}{\mathcal{G}} = -1\,.              \label{Eq:grp122}
\end{equation}
The gravitational potential energy of the $n=0$ polytropes is represented in
Fig.\,\ref{UniSigGBEn}. The role of the cosmological constant is illustrated
by the sequence of lines constructed for appropriately chosen values of the
cosmological parameter~$\lambda$.

In polytropes with $n>0$, the binding energy $\mathcal{B}$ differs from the
gravitational binding energy $(-\mathcal{G})$, as some of the work of the
gravitational field is converted into the kinetic energy of microscopic motion
in these polytropes. The role of the cosmological parameter $\lambda \neq 0$
in the effect of gravitational binding will be discussed in the next section.

\endinput

%%
%% End of file `x_poluni.tex'.
